\scanpage{32}
% Pencil archive foliation: 31.
plus bas).]
\margin{\deleted{bon pour}}
\deleted{$\mathcal{M}^i(X)$ est ``engendré'' par les
$R^if_*$ L'appartenance à $\mathcal{M}^i(\unsure{k(x)|X})$ de
$M|_x$ implique \unsure{que}}
Compatible aux
$f^*$, et \deleted{l'appart.}
\[
M|_x\in\operatorname{Ob}\mathcal{M}^{+i}(k(x))
\Longrightarrow
\exists U\ni x\ \text{voisinage de $\bar{x}$, tel que }
M|_U\in\operatorname{Ob}\mathcal{M}^{+i}(U).
\]
\margin{\unsure{quasi-limite} Compatibilité aux $\underset{\longleftarrow}{\lim}$ de préschémas \ldots}
Soit $f:X\unsure{-|\longrightarrow}Y$,
alors
\[
R^jf_!:\mathcal{M}^{+i}(X)\longrightarrow \mathcal{M}^{+i+j}(Y);
\]
de plus, $\mathcal{M}^{+i}(k)$ ($k$ un corps alg. clos)
est ``engendré'' par les \inserted{sous-quotients}
$R^jf_!(\deleted{\mathbb{Z}}\mathbf1_X)$ pour
$X$ \unsure{disons} projectif lisse de dim $\leq j$ sur $k$
[au sens que tout $M\in\operatorname{Ob}\mathcal{M}^i(k)$ a une
filtration dont les facteurs sont isomorphes
à de tels sous-quotients].
\margin{indices en bas}
\[
\deleted{\mathcal{M}^i(X)\otimes\mathbb{Z}(-1)}
\]
\[
\left\{
\begin{array}{l}
\mathcal{M}_i^+(X)\otimes \mathcal{M}_j^+(X)\subset \mathcal{M}_{i+j}^+(X)\\
\mathbb{Z}(-1)\in\operatorname{Ob}\mathcal{M}_2^+(X)
\end{array}\right.
\]
d'où
\[
\mathcal{M}^{+i}(X)\otimes\mathbb{Z}(-j)\subset \mathcal{M}^{+i+2j}(X)
\qquad\text{pour }j\geq0
\]
